\documentclass[10pt]{article}
\usepackage{template}
\usepackage{amsmath,amssymb,microtype,tikz,graphicx,xurl}
\usetikzlibrary{arrows.meta,positioning,calc}
\setlength{\emergencystretch}{3em}
\raggedbottom
\SetKeys{date=2026-10-29,vol=016,title={每日有趣新闻｜时政、科学与数学},author={ChatGPT 编辑部},note={\small 采编窗口：2026年10月27—29日；截稿：北京时间10月29日上午。政治报道区分政府主张、媒体核实与本刊推断；数学预印本及其形式化声明不等同于独立复核。}}
\renewcommand{\newsitem}[3]{\par\medskip{\large\bfseries\color{newsblue}#1\par}{\footnotesize\color{newsgray}#2\par}\smallskip\color{black}#3\par\medskip{\color{newsgray}\rule{\linewidth}{0.25pt}}\par}
\newcommand{\newslead}[3]{\par\medskip{\LARGE\bfseries\color{newsblue}#1\par}\vspace{1mm}{\small\color{newsgray}#2\par}\smallskip\color{black}#3\par\medskip{\color{newsgray}\rule{\linewidth}{0.45pt}}\par}
\newcommand{\newssub}[1]{\par\smallskip\noindent{\bfseries\color{newsblue}#1}\quad}
\begin{document}
\begin{documents}
\raggedcolumns

\newslead{联合国再促解除对古巴禁运：赞成票减少，外交阵线变化}{国际政治深度｜联大10月28日表决；对照联合国2025年正式记录、路透社与美联社}{
联合国大会10月28日再次通过要求美国终止对古巴经济、商业和金融禁运的决议。\textbf{赞成130票、反对17票、弃权25票}。从表决结果看，赞成者仍显著多于反对者；从年度变化看，这却是一次值得细究的外交转向。联合国正式记录显示，2025年的相应票数为165票赞成、7票反对、12票弃权。两年的“赞成”减少35票，“反对”增加10票，“弃权”增加13票；剩余国家未投票，不能把全部差额都解释成转投反对。\cite{ReutersCuba,APCuba,UN2025}
\begin{center}
\begin{tikzpicture}[x=0.030cm,y=0.72cm, font=\scriptsize]
\node[anchor=west,font=\bfseries] at (0,2.4) {联大古巴禁运决议投票变化（票数）};
\foreach \y/\yr/\yes/\no/\abs in {1.5/2025/165/7/12,0.4/2026/130/17/25}{
  \node[anchor=east] at (-2,\y+.15) {\yr};
  \fill[newsblue!80] (0,\y) rectangle (\yes,\y+.36);
  \fill[orange!75] (\yes,\y) rectangle ({\yes+\no},\y+.36);
  \fill[gray!55] ({\yes+\no},\y) rectangle ({\yes+\no+\abs},\y+.36);
  \node[anchor=west] at (0,\y-.18) {赞成\yes，反对\no，弃权\abs};
}
\node[anchor=west,newsblue] at (0,-.28) {\rule{1.8mm}{1.8mm}\,赞成};
\node[anchor=west] at (50,-.28) {\textcolor{orange}{\rule{1.8mm}{1.8mm}}\,反对};
\node[anchor=west] at (99,-.28) {\textcolor{gray}{\rule{1.8mm}{1.8mm}}\,弃权};
\end{tikzpicture}
\par\vspace{1mm}{\scriptsize\color{newsgray}图1｜两年已记录的三类投票，不含未投票国家；横轴长度按票数比例。2025年联合国记录，2026年路透社、美联社10月28日报道；本刊据数据重绘。}\end{center}\color{black}

\newssub{一份没有直接强制力的决议，为何年年投票？}联合国大会决议通常不具备安全理事会根据《联合国宪章》第七章采取强制措施那样的法律效力。美国的禁运政策植根于美国国内立法及行政措施，联大多数票不会自动解除美国国会制定的限制。因此，这次表决首先衡量的是\textbf{各国愿意公开承担何种外交立场}，而不是一个马上生效的贸易协定。自1992年以来，联大已多年重复这一议题，2020年因疫情未按通常安排表决。长期多数并未直接改变禁运制度，正说明政治谴责与国内法律之间存在制度距离。\cite{ReutersCuba,APCuba}

\newssub{变动发生在哪里？}美联社列举德国、意大利和拉脱维亚等国今年改投弃权；波黑、哥伦比亚、哥斯达黎加、厄瓜多尔、立陶宛、摩洛哥、北马其顿、巴拉圭和斯洛文尼亚等国列入反对阵营。美国驻联合国代表迈克·沃尔兹公开将这种变化视为美国外交游说的成果。德国、意大利等国的投票变化，意味着传统上反对禁运的一些伙伴不再愿意以相同方式与美国划清界限。但一国弃权的具体动机，可能涉及双边谈判、联盟协调、国内政治及对古巴政府的态度；仅凭一张票不能断言每个国家接受了美国全部论述。\cite{APCuba,ReutersCuba}

\newssub{双方在争论什么？}美国政府主张，对古巴施压有助于推动其政治制度变化，并批评古巴政府的治理、人权及政治自由状况。古巴政府则指责美国制裁及燃料限制造成严重民生伤害，将其视为侵犯主权的集体惩罚。路透社采访哈瓦那居民，记录到能源短缺、停电、食品和饮水供应困难等现实困境；但这些观察本身无法把每一项经济损失精确分摊到制裁、古巴国内经济政策、基础设施老化及其他因素。对经济伤害的归因需要贸易、能源、财政与民生的连续数据，不能只引用任一政府的宣传口径。\cite{ReutersCuba,APCuba}

\newssub{会场争执与外交事实须分开。}表决前后，美古代表在联合国发生激烈口角。美国方面指称古巴一名外交官袭击美方人员，并宣布将其驱逐；古巴方面则向联合国秘书长提出针对美方人员的投诉。双方相互指责已被媒体记录，\textbf{但谁对具体肢体冲突负有责任，不能仅凭当事国声明裁定}。把外交争执写成已经获得独立调查确认的袭击事实，会掩盖证据等级的区别。\cite{ReutersCuba,APCuba}

\newssub{如何判断这次表决的意义？}有两件事同时成立：其一，国际社会多数成员依旧不支持美国对古巴长期实施禁运；其二，美国在今年显著削弱了过去的压倒性多数。前者说明禁运在国际政治上仍缺乏广泛认同，后者说明这种认同并非不会改变。是否解除制裁取决于美国国内政策，是否改善古巴民生则还取决于古巴自身治理和国际贸易条件。对普通古巴居民来说，比会场胜负更重要的是燃料、药品、电力和收入能否实际改善。把联合国的外交象征、美国的政策工具与古巴人的生活放在三个不同层次，才能理解这条新闻。\cite{ReutersCuba,APCuba}
}

\newslead{火星尘埃比预想的更细：毅力号测到百纳米颗粒，宇航服该如何防护？}{航天与行星科学深度｜NASA/JPL 10月26日发布；PNAS研究10月26日在线发表，28日媒体报道}{
火星探测器的主要任务通常是寻找古老水环境、分析岩石，或者为样本返回做准备。NASA的“毅力号”却利用仪器校准靶和随车携带的宇航服材料样片，回答了一个关乎未来人类登陆的工程问题：\textbf{火星尘埃究竟有多细，它会怎样附着在人的装备上？}Scott J. VanBommel等人在《美国国家科学院院刊》（PNAS）发表的研究表明，火星尘粒不仅有此前可测的约3微米级颗粒，还包含小于2.5微米的细颗粒，以及尺度约100至200纳米的超细颗粒。这里的“发现”是对已有火星表面尘埃的更细尺度测量，并不是突然发现了一种全新的火星物质。\cite{PNASDust,NASADust}
\begin{center}
\begin{tikzpicture}[x=0.75cm,y=0.75cm, font=\scriptsize,>=Stealth]
\node[font=\bfseries,anchor=west] at (0,3.3) {火星尘粒：研究测得的尺度};
\draw[->,newsgray] (0.3,.4) -- (8.8,.4) node[right] {粒径};
\foreach \x/\lbl in {1/0.1,3/1,5/10,7/100}{\draw (\x,.3)--(\x,.5);\node[below] at (\x,.3) {\lbl $\mu$m};}
\fill[orange!70] (1,1) circle (.09); \node[anchor=west] at (1.35,1) {超细尘：约$0.1\,\mu$m};
\fill[newsblue!65] (3.8,1.8) circle (.23); \node[anchor=west] at (4.2,1.8) {细尘：低于$2.5\,\mu$m};
\fill[gray!55] (7.0,2.5) circle (.40); \node[anchor=east] at (6.5,2.5) {人类头发约$70\,\mu$m};
\end{tikzpicture}
\par\vspace{1mm}{\scriptsize\color{newsgray}图2｜对数位置仅用于数量级比较，圆面积不按比例。数据：NASA/JPL 2026-10-26新闻稿及VanBommel等PNAS论文；本刊原创示意，非显微照片。}\end{center}\color{black}

\newssub{测量是怎样完成的？}毅力号携带的PIXL（行星X射线岩石化学仪）主要通过X射线荧光分析目标的元素组成。探测器上有供仪器校准的材料，包括聚四氟乙烯（Teflon）与玄武质玻璃；另有供SHERLOC仪器使用的Vectran纤维和Teflon样片。随着火星风吹起尘土，这些样片成了天然的长期沉降试验台。研究团队利用PIXL的观测和模型分析材料表面尘粒的尺度及化学信号。NASA指出，Teflon本身不含PIXL可直接检测的目标元素，因此在其表面观测到的相应元素信号可以用于分析沉积尘埃。\cite{NASADust,PNASDust}

\newssub{为何“百纳米”是个重要尺度？}1微米是千分之一毫米，100纳米等于0.1微米。NASA的新闻稿以约70微米的人类头发作比较，0.1微米约为其七百分之一。空气中悬浮的颗粒越小，越可能绕过鼻腔和上呼吸道的部分防御机制，进入更深的肺部；但\textbf{粒径小并不自动等于毒性大}。健康风险还取决于吸入剂量、暴露时间、化学组成、颗粒表面活性和人体防护条件。研究发现的富含铁、硫、磷等元素的硅酸盐尘埃，以及火星环境中已知的某些含氯化合物，都要求工程师谨慎评估。没有宇航员在火星实际暴露的临床数据，不能据此推算具体患病概率。\cite{PNASDust}

\newssub{材料差异比“哪种更耐磨”复杂。}NASA报道，Vectran样片更容易截留较大的颗粒，部分集中在织物背面；Teflon附着的尘埃总体较少，但留下的颗粒可能更细。静电吸附和材料表面结构可能共同决定这一差异。研究没有证明这些样片已遭严重材料损伤，也没有给出“某种宇航服材料一定安全”的最终结论。要区分\textbf{附着量、粒径分布、毒性和材料损耗}四个指标：一种材料不易粘尘，仍可能让极细颗粒更难清除；另一种材料捕获较多颗粒，也不意味着它一定更易破损。\cite{NASADust}

\newssub{火星与月球的灰尘不一样。}阿波罗宇航员曾报告月尘磨损面罩、密封件和外层织物。月球几乎没有能长期风化颗粒的大气，微陨石撞击形成的碎屑常有锋利棱角。火星则有风和持续的颗粒碰撞，许多尘粒被磨得更圆，机械磨蚀性未必与月尘相同；火星尘埃更突出的挑战可能是细颗粒的呼吸风险、化学性质以及持续进入居住舱的可能性。未来的“宇航服接口”（suitport）和双层气闸舱，正是为了减少宇航服从室外把尘埃带进生活空间。\cite{NASADust}

\newssub{什么还没有得到证明？}这项研究直接支持的是“存在此前难以测到的细尘，并在不同材料上表现出不同附着特征”。从这个结果走到载人火星任务的实际风险评估，还需要长期暴露实验、尘埃毒理、静电清洁技术、密封系统测试，以及不同地形和季节的重复观测。NASA强调这是一个需要设计解决的风险，不是火星不可居住的宣判。科学新闻的有趣之处恰在这里：一块原本用于校准的材料，能够在多年后变成研究人类未来生活环境的样本。\cite{PNASDust,NASADust}
}

\newsitem{美联储维持3.75\%—4.00\%利率区间：两名委员主张继续加息}{经济政策重点｜美联储10月28日14:00美东时间正式声明；美联社核对}{
美国联邦公开市场委员会（FOMC）在10月27日至28日的会议后决定，维持联邦基金利率目标区间在3.75\%至4.00\%。\textbf{表决为10比2}：Lisa D. Cook和Beth M. Hammack两名委员主张当次再提高25个基点。官方声明称经济活动仍在较强速度扩张、就业增幅与劳动力增长大致匹配，同时明确通胀仍然偏高。9月的上次会议曾加息25个基点，因此本次是\textbf{暂停继续上调}，而不是降息或转向宽松。\cite{FedStatement,APFed}

联邦基金利率是美国银行间隔夜资金市场的重要政策利率，并非所有住房贷款、企业贷款都直接按这一数值定价。它通过银行融资成本、债券收益率、资产价格和汇率等渠道影响经济。美联储的长期通胀目标为2\%；美联社称近期通胀数据仍高于这一水平。政策利率升高通常抑制信贷需求，但传导存在时滞，能源价格和供给冲击也未必能靠提高利率迅速解决。\cite{FedStatement,APFed}

值得注意的是，声明的措辞强调需要积累更多信息，以确定能让通胀及时回到目标的适当政策立场。市场人士押注12月会议继续加息，是市场定价与经济预测，不是FOMC预先承诺。距美国11月中期选举只剩数日，主席Kevin Warsh公开否认政治因素左右决定；这一表态可以准确报道，却不等于记者已经获得了对每位委员全部动机的独立验证。\cite{APFed}

本次决议最能说明的，是美联储在通胀与经济活动之间继续权衡：不加息并不代表认为通胀问题已经结束；两票反对也不意味着委员会整体支持立即收紧。未来需观察后续通胀、就业、信贷和12月实际表决，而非把一场记者会的措辞当作政策保证。
}

\newsitem{有机分子涂层也能成为量子芯片的一部分？从医学成像到量子器件的意外路线}{物理与化学重点｜Nature 10月28日论文；路透社采访德国NVision}{
德国初创公司NVision的研究团队10月28日在《自然》发表《单分子自旋—光子接口》（A single-molecule spin--photon interface）。研究者将具有三重态基态的卡宾分子嵌入结构匹配的晶体，展示了对\textbf{单个分子}的光学寻址、光学读出磁共振和相干操纵。实验在4.5开尔文条件下得到约2.2毫秒的动态去耦自旋相干时间，以及约21毫秒的自旋弛豫时间。路透社指出，公司早期的自旋技术源自磁共振成像研究；本次工作则把分子自旋与光学接口结合，探索可与光子集成电路相连接的量子器件。\cite{NatureNV,ReutersNV}

量子比特不是“把普通比特做得更小”。普通数字电路用可区分的高低电平编码0和1；量子器件则要维持并控制叠加态及量子相干，还要能可靠初始化、操纵、读取，并控制噪声和误差。电子自旋是候选物理载体之一。OLED显示产业已经拥有成熟的有机薄膜沉积技术，因此如果某类自旋分子能兼容这套工艺，就可能绕开部分传统量子硬件制造难题。\cite{ReutersNV,NatureNV}

路透社称，团队先用AI筛选数百种候选分子，再实际合成约40种，最终选择适合薄膜沉积的材料。这是一条\textbf{计算筛选—化学合成—器件验证}的工作链，不能将其简写成“AI直接发明量子计算机”。公司宣布与美国芯片企业HyperLight合作制造首批芯片，并计划未来开发完整系统；但量产良率、可扩展量子比特数量、门保真度、纠错开销和实际算力尚待公开检验。\cite{ReutersNV}

还需区分经同行评议发表的具体实验研究与企业的长期商业目标。Nature论文经同行评议，证明的是其报告的单分子自旋—光子接口及低温相干操纵，并未证明已交付通用容错量子计算机，更不能将4.5开尔文下的实验结果直接外推到室温芯片。相较那些把“量子”直接等同于计算优势的宣传，这条新闻更值得关注的，是一种源自医学成像的分子技术如何跨越学科边界，变成制造量子器件的新候选路线。
}

\newsitem{数学拾趣｜一个元素同时“本原”又“完全正规”：30年有限域猜想的新证明主张}{数论与Lean/mathlib重点｜Florian Hellwig；arXiv:2610.36534，10月27日提交}{
在有限域$\mathbb F_{q^n}$中，$q$是素数幂，$n$是扩张次数。非零元素组成一个有$q^n-1$个元素的乘法循环群；若$\alpha$的幂能遍历全部非零元素，就称它为\textbf{本原元素}。另一方面，如果$\alpha,\alpha^q,\ldots,\alpha^{q^{n-1}}$在$\mathbb F_q$上构成向量空间的一组基，就称$\alpha$为\textbf{正规元素}。本原性关注乘法结构，正规性关注加法及线性代数结构，它们并不是同一回事。\cite{HellwigPaper}

更苛刻的要求是\textbf{完全正规}：对每个介于$\mathbb F_q$和$\mathbb F_{q^n}$之间的中间域$K$，同一个$\alpha$都必须在扩张$\mathbb F_{q^n}/K$中是正规元素。Morgan与Mullen在1996年猜想：每个有限域扩张都存在既本原又完全正规的元素。先前已有部分参数范围的定理，但把所有素数幂$q$与所有正整数$n$统一处理，需要同时控制多层线性条件和乘法生成条件。\cite{HellwigPaper,HellwigCode}
\begin{center}
\begin{tikzpicture}[font=\scriptsize, node distance=8mm and 16mm, >=Stealth,
 every node/.style={draw=newsblue!70,rounded corners=2pt,inner sep=4pt}]
\node (top) {$\mathbb F_{2^6}$};
\node (midl) [below left=9mm and 13mm of top] {$\mathbb F_{2^2}$};
\node (midr) [below right=9mm and 13mm of top] {$\mathbb F_{2^3}$};
\node (bottom) [below right=9mm and 13mm of midl] {$\mathbb F_2$};
\draw[->] (bottom)--(midl);\draw[->] (bottom)--(midr);
\draw[->] (midl)--(top);\draw[->] (midr)--(top);
\end{tikzpicture}
\par\vspace{1mm}{\scriptsize\color{newsgray}图3｜有限域$\mathbb F_{2^6}$的部分中间域；箭头表示包含关系。完全正规性要求对每个中间域分别成立。本刊依据有限域结构绘制。}\end{center}\color{black}

一个容易动手验证的小例子是$\mathbb F_4=\mathbb F_2[\omega]/(\omega^2+\omega+1)$。$\omega^3=1$且$\omega\ne1$，故它生成三个非零元素；$\omega$与$\omega^2=\omega+1$又在$\mathbb F_2$上线性无关，构成正规基。这个例子展示两种条件可以共存，但显然不能代替对全部$q,n$的证明。

Hellwig的新预印本使用部分拉格朗日预解式把完全正规性拆分成可控制的分量，结合特征和估计与有限计算，宣称证明了1996年的猜想，并得到完全正规元素计数下界及与本原性的渐近独立性。论文称对$n>q$的参数处理，除70对$(q,n)$需计算机给出明确元素外，可通过统一准则解决；$n\le q$则引用已发表的Garefalakis—Kapetanakis定理。\cite{HellwigPaper}

\newssub{Lean究竟验证了什么？}作者公开了Lean 4/mathlib代码仓库，README宣称`lake build`成功、无`sorry`，且列出了主定理依赖的公理清单。关键边界非常具体：$n\le q$的已发表定理和Rosser—Schoenfeld的一项估计被作为外部公理输入；部分大型有限计算通过`Lean.ofReduceBool`及`Lean.trustCompiler`执行，信任基础不仅包括Lean内核，也包括相应编译器、解释器和底层算术实现。作者提供了25项检查器与审计脚本，但\textbf{本刊没有在独立环境中下载依赖、运行全部构建、审计形式化定义与外部公理}。因此准确表述是“作者提交了带有明确可信基础说明的形式化证明”，而不是“本刊已证实30年猜想彻底解决”。\cite{HellwigCode}

这份研究尤其值得关注，因为它把数学论文中容易含糊的一句“计算机已验证”，拆成了可追问的层次：数学陈述是否正确翻译成形式语言、依赖哪些外部定理、计算结果由谁检查、证明内核究竟负责到哪一步。即使将来发现需要修订，公开的可信基础与可复现代码仍有独立的方法论价值。
}

\newsitem{默茨在东京提议贸易伙伴联盟：德日希望在大国关税竞争中增加选择}{外交与经济重点｜10月28日德国总理访日；路透社、双方公开表态}{
德国总理默茨10月28日在东京呼吁日本支持由“志同道合国家”组成的贸易合作联盟，以协调应对关税冲突及经济胁迫。日本首相高市早苗欢迎德国扩大与亚太伙伴合作；双方联合表态强调以规则为基础的贸易，并表示将探索防务产业合作的可能性。默茨还提出加强欧盟与《全面与进步跨太平洋伙伴关系协定》（CPTPP）成员之间的联系。\cite{ReutersMerz}

这项倡议目前是\textbf{外交提议}，不是已经成立、具备章程和执行机构的新国际组织。德国的现实动机包括出口市场变化、对中国供应链的依赖，以及在美国关税政策下争取更多谈判空间。默茨提及日本在2010年稀土供应争端后扩大储备、寻找替代来源和发展回收技术的经验。稀土、半导体材料和能源并非简单靠一个自由贸易协议就能获得供应保障；库存、替代工艺、回收能力和长期合同都属于不同层次的风险管理。\cite{ReutersMerz}

中国外交方面呼吁德国和欧盟对经贸关系保持客观理性，并表示愿维护供应链稳定。这是中方政策立场，不应写成已接受德国提议。德国政府强调对华“去风险”，也不能自动等同于全面经济脱钩：降低单一供应来源的脆弱性与停止贸易，是程度差异极大的两种政策。\cite{ReutersMerz}

德日同时讨论潜在的联合战机研发，但目前只是探索可能性，不代表项目已经立项、签署合同或确定技术路线。未来最值得追踪的不是会谈合影，而是参与国名单、具体贸易争端协调机制、供应链项目投资以及这些合作是否与现有欧盟和CPTPP规则兼容。
}

\newsitem{宇树科技三季度营收增长70\%，净利润却下降：机器人商业化要看什么？}{中国科技产业｜10月28日披露三季度财务数据；路透社报道}{
宇树科技10月28日披露上市后的首份季度业绩。第三季度营收为\textbf{6.65亿元人民币}，同比增长约70\%；归属于股东的净利润约\textbf{1.127亿元}，同比下降17.4\%。研发费用约1.27亿元，同比增加249.6\%，相当于当季收入的19.1\%。这些数据描述的是不同财务指标：营收增长不保证利润同步增长，研发费用增加也不能直接推断研发失败。\cite{ReutersUnitree}
\begin{center}
\begin{tikzpicture}[x=0.045cm,y=0.64cm,font=\scriptsize]
\node[anchor=west,font=\bfseries] at (0,2.6) {宇树科技2026年第三季度：同比变动};
\draw[gray!60] (0,.3)--(0,2.35);
\foreach \y/\label/\v in {2/营收/70,1.3/研发支出/249.6,.6/归母净利润/-17.4}{
 \node[anchor=east] at (-4,\y+.12) {\label};
 \ifdim\v pt>0pt \fill[newsblue!75] (0,\y) rectangle ({\v/3},\y+.26);\else \fill[orange!80] ({\v/3},\y) rectangle (0,\y+.26);\fi
 \node[anchor=west] at ({max(0,\v/3)+3},\y+.12) {\v\%};
}
\end{tikzpicture}
\par\vspace{1mm}{\scriptsize\color{newsgray}图4｜同比变化率，负值向左；各指标不是同一金额单位。路透社据宇树上海证券交易所财报2026-10-28报道，本刊重绘。}\end{center}\color{black}

前九个月，公司营收约18.2亿元，同比增长55.7\%；归母净利润约3.867亿元，同比增长270.5\%。公司解释，前九个月利润增幅较大，部分受到上一年一次性股份支付费用形成的低比较基数影响。因此，\textbf{单季度利润同比下降与前三季度累计利润同比大增可以同时成立}，并不矛盾。评价机器人企业，还需要分清四足机器人与人形机器人的销售构成、硬件毛利、软件收入、交付后的维护成本和实际应用场景。\cite{ReutersUnitree}

资本市场将宇树视为中国具身智能产业的代表，但舞蹈、武术或奔跑视频主要展示运动控制能力，不足以证明机器人在复杂工厂、家庭护理或物流环境中的长期经济性。此前股价大幅波动，意味着市场对未来增长的预期变化剧烈。将企业财报与技术演示放在一起观察，比只凭股价涨跌判断中国机器人技术的进退更可靠。\cite{ReutersUnitree}
}


\newsitem{伊朗停战斡旋遇阻：新的建议并不等于达成协议，战争账单已达471亿美元}{中东战争与外交重点｜路透社10月28日；美联社援引美国中央司令部截至10月16日估算}{
10月28日，卡塔尔斡旋方就结束伊朗战争提出新一轮建议。路透社援引一名伊朗官员称，德黑兰收到的方案包含其认为不可接受的内容；另一消息来源则描述了若干建立互信的步骤。两种说法并不必然矛盾：调解者可以提出逐步降级方案，而交战方仍可能认为关键条件不够。\textbf{有建议、有接触、有表态，都不等于已签署停火协议}。战事还与也门胡塞武装袭击沙特炼油设施、国际能源运输安全等议题相互交织，增加了地区谈判的复杂性。\cite{ReutersIran}

与此同时，美联社10月28日报道，美国中央司令部向国会提供的最新估算显示，\textbf{截至10月16日，美国参与伊朗战争的成本约为471亿美元}，比9月初公布的约436亿美元增加35亿美元。这是美国军方按特定统计口径估计的战争支出，并不是全球战争总经济损失，也不包括伊朗及周边国家的全部人员、基础设施和社会成本。不能将其直接理解成国会当天又通过了471亿美元新拨款。\cite{APIranCost}

战争成本至少有三个层次。直接财政支出包括武器弹药、舰机行动、部署、维护和人员保障；间接经济成本包括能源价格、航运保险、供应链延误与通胀；长期成本还包括伤亡者医疗照护、基础设施重建、投资流失和政治风险。第一类可以由政府会计估算，后两类则需要更复杂的反事实模型：若没有战争，油价和投资本来会怎样变化？这类问题不可能从一张账单直接读出答案。

美国政府以国家安全与地区威慑解释军事行动，伊朗方面则强调主权与安全要求。对这些政治立场应作来源归属，不能把任何一方的战争目标当作客观已证实的正当性结论。衡量斡旋进展，更可靠的指标是双方是否就攻击停止时间、监督机制、被拘押人员、核活动、制裁与后续谈判议程形成可执行文本，而不是仅看新闻标题里是否出现“和平方案”。\cite{ReutersIran,APIranCost}
}

\newsitem{美国再制裁涉嫌为朝鲜筹资的跨国IT工作网络：招聘身份核验成为安全问题}{网络安全与国际政策｜美国OFAC 10月28日制裁记录；路透社报道}{
美国财政部海外资产控制办公室（OFAC）10月28日更新朝鲜相关制裁名单。路透社报道称，制裁针对一个被美国指控利用虚假身份为境外企业提供IT服务、并将所得收入输送给朝鲜的网络。\textbf{制裁指定是美国行政机关的法律行为；有关个人和机构的违法指控仍应归属于美国政府，不能自动写成国际司法机关已判决的事实}。\cite{OFAC2026}

这类活动与传统“黑客直接攻破服务器”并不完全相同。它可能从招聘流程进入：求职者提供经过包装的履历、身份资料和远程办公地址，通过代理网络、转包或中介完成工作。企业即使没有遭遇恶意软件，也可能在工资支付、客户资料、源代码访问和制裁合规方面承担风险。美国财政部过去公布的案例还提及部分人员利用虚假身份取得软件开发合同；这说明相关风险有历史背景，但不意味着所有海外远程开发者都应受到怀疑。\cite{OFACBackground}

企业可采取的应对不只是简单屏蔽某些国家的IP地址，而是对敏感岗位开展与风险相称的身份及付款账户核验、权限最小化、代码审查和供应商尽调，同时避免把国籍或远程办公本身当作违法证据。安全措施若以过度监控或歧视为代价，也可能制造新的合规问题。\cite{OFACBackground}
}


\newsitem{数学短论｜布尔函数的重量谱：一项Reed--Muller码猜想遭到反例挑战}{编码理论与组合数学｜Yueying Lou、Qichun Wang；arXiv:2610.36687，10月27日提交}{
布尔函数是从$\{0,1\}^m$到$\{0,1\}$的映射。最简单的例子是“异或”函数$f(x,y)=x+y\pmod2$，在四组输入中有两组输出1。\textbf{重量}就是函数输出1的输入个数，因此这个函数的重量为2。把所有次数不超过$r$的$m$元布尔多项式收集起来，得到经典的Reed--Muller码$\mathrm{RM}(r,m)$；所有码字可能出现的重量组成\textbf{重量谱}。这个集合不像仅计算码字总数那样容易：不同代数形式会出现相同重量，而有些看似可能的重量其实无法实现。\cite{LouWang}

为什么研究重量谱？编码理论关心数据传输时哪些错误模式容易混淆。两个二进制码字相差的位数称为汉明距离；线性码的非零码字重量与距离分布紧密相关。重量谱不仅有助于研究纠错能力，还与布尔函数的代数结构、密码学中函数的非线性性质和有限几何有联系。特别是当$r$接近$m$时，码字非常多，要判定某个整数究竟能否成为重量，往往需要比直接枚举更精细的方法。

Carlet与Solé在2023年研究$\mathrm{RM}(m-c,m)$的重量谱，并提出对较一般参数的猜想。此后$c=5$、$c=6$的情形取得进展，给了猜想一定支持。Lou与Wang这篇新预印本声称，针对每个$m\ge14$，确定$\mathrm{RM}(m-7,m)$的精确重量谱；尤其在$\mathrm{RM}(7,14)$中，他们使用仿射几何与傅里叶分析排除了猜想所预期的八个重量，再通过几何方法推广到更高维。若证明成立，就构成对原猜想的反例。\cite{LouWang}

这里需要注意两个逻辑区别。首先，\textbf{否定一个重量谱猜想}不等于否定Reed--Muller码的基本理论，也不意味着这些码失去了纠错用途。它只是说明某个关于可实现重量的预测过强。其次，作者声称“确定了精确谱”，比“找到了一个反例”更强；完整核验不仅要确认被排除的重量确实不可能出现，还要确认其余被列入的重量都有构造或证明。论文目前是arXiv预印本，本刊尚未独立检查全部证明，亦未发现其主结果已获得公开的Lean内核验证。\cite{LouWang}

这类研究还有一层方法论趣味：一个猜想可以在多个低维实例上成立，却在更高维遇到结构性障碍。有限样例能提供直觉，但不能代替统一定理。与本期有限域的形式化证明主张放在一起看，恰好呈现了数学研究的两种方向：一边努力把存在性命题证明到所有参数，一边用细致的结构分析寻找旧猜想的边界。
}

\par\medskip{\large\bfseries\color{newsblue}参考资料与编辑说明}\par
{\scriptsize 本期主要检索2026年10月28日新闻与原始材料，另纳入10月26—27日发表、28日获广泛报道的科学论文。四张插图均为本刊原创TikZ图表/示意，不使用未经授权的路透社或美联社现场摄影；数字均标注日期及口径。论文的同行评议、代码声明与本刊独立核验严格区分。}\par
\begingroup\fontsize{6}{7.1}\selectfont
\renewcommand{\refname}{参考文献}
\bibliographystyle{plainurl}\bibliography{refs}
\endgroup
\end{documents}
\end{document}